\section{Introducción}

Las plataformas de música registran cada reproducción como un evento: quién escuchó, qué canción, de qué artista y en qué momento. Con millones de estos eventos se puede estudiar cómo se reparte la escucha en el tiempo, qué artistas concentran la atención de los usuarios y qué relaciones aparecen entre gustos musicales. Este proyecto, al que llamamos \emph{Music DNA}, construye una arquitectura de datos para hacer ese análisis sobre datos públicos de ListenBrainz \citep{listenbrainzdata}.

ListenBrainz es un proyecto de la fundación MetaBrainz que recopila las escuchas que envían sus usuarios y publica volcados periódicos en formato JSON Lines bajo licencia CC0. Trabajamos con un volcado incremental concreto, el número 2746 del 27 de octubre de 2026, que contiene 4\,855\,809 escuchas de 26\,394 usuarios. Fijar un volcado hace que todo el grupo trabaje sobre los mismos datos y que los resultados se puedan reproducir.

El objetivo del proyecto es almacenar, procesar y analizar esas escuchas con una arquitectura multimodelo: MongoDB \citep{mongodb} guarda los documentos originales, Apache Cassandra \citep{cassandra} guarda la serie temporal y las métricas agregadas por fecha, y Neo4j \citep{neo4j} modela las relaciones entre usuarios, canciones y artistas. Apache Airflow \citep{airflow} orquesta el flujo completo, desde la descarga del volcado hasta el cálculo de los indicadores, y deja los resultados en un formato pensado para un dashboard en Streamlit \citep{streamlit}.

El informe sigue el orden de la arquitectura. La sección~\ref{cap:problema} plantea el problema y por qué una sola base de datos no alcanza. La sección~\ref{cap:arquitectura} describe el stack completo y justifica qué dato va en cada base. La sección~\ref{cap:dataset} presenta el dataset y su validación. Las secciones~\ref{cap:mongodb}, \ref{cap:cassandra} y~\ref{cap:neo4j} detallan el modelo de datos, la carga, el CRUD y las consultas de cada base. La sección~\ref{cap:airflow} explica el DAG de Airflow, la sección~\ref{cap:kpis} define los cinco indicadores y la sección~\ref{cap:dashboard} presenta el diseño del dashboard. La sección~\ref{cap:conclusiones} cierra con lo aprendido.
